\documentclass[letterpaper,11pt]{article}

\usepackage[margin=0.75in]{geometry}
\usepackage[hidelinks]{hyperref}
\usepackage{titlesec}
\usepackage{enumitem}
\usepackage{fancyhdr}
\usepackage{lastpage}
\usepackage[usenames,dvipsnames]{xcolor}
\usepackage[style=ieee,url=false,doi=true,maxbibnames=99,sorting=ydnt,dashed=false]{biblatex}
\bibliography{papers}

% ATS-friendly: minimal formatting, no graphics, no columns, standard fonts
\pagestyle{fancy}
\fancyhf{}
\renewcommand{\headrulewidth}{0pt}
\fancyfoot[C]{\small Page \thepage\ of \pageref{LastPage}}

\definecolor{accent}{HTML}{1a1a2e}
\titleformat{\section}{\large\bfseries\color{accent}\uppercase}{}{0em}{}[\titlerule]
\titlespacing*{\section}{0pt}{12pt}{6pt}

\setlist[itemize]{leftmargin=*, itemsep=1pt, parsep=0pt, topsep=2pt}
\setlist[enumerate]{leftmargin=*, itemsep=1pt, parsep=0pt, topsep=2pt}

% Bold author name in bibliography
\newcommand*{\boldname}[3]{%
  \def\lastname{#1}\def\firstname{#2}\def\firstinit{#3}}
\boldname{Guggilam}{Sreelekha}{}
\renewcommand{\mkbibnamegiven}[1]{%
  \ifboolexpr{ ( test {\ifdefequal{\firstname}{\namepartgiven}} or test {\ifdefequal{\firstinit}{\namepartgiven}} ) and test {\ifdefequal{\lastname}{\namepartfamily}} }
  {\mkbibbold{#1}}{#1}}
\renewcommand{\mkbibnamefamily}[1]{%
  \ifboolexpr{ ( test {\ifdefequal{\firstname}{\namepartgiven}} or test {\ifdefequal{\firstinit}{\namepartgiven}} ) and test {\ifdefequal{\lastname}{\namepartfamily}} }
  {\mkbibbold{#1}}{#1}}

\begin{document}

% ============================================================================
% HEADER -- ATS reads top-down, single column, no tables
% ============================================================================
\begin{center}
{\LARGE \textbf{Sreelekha Guggilam, Ph.D.}}\\[4pt]
Assistant Professor in Data Science\\
Department of Mathematics and Statistics, Texas A\&M University--Corpus Christi\\[4pt]
Email: \href{mailto:sreelekha.guggilam@tamucc.edu}{sreelekha.guggilam@tamucc.edu}
\enspace$\vert$\enspace Phone: +1 361-825-3485
\enspace$\vert$\enspace LinkedIn: \href{https://www.linkedin.com/in/sreelekhaguggilam}{sreelekhaguggilam}\\
GitHub: \href{https://github.com/Datalakes9}{Datalakes9}
\enspace$\vert$\enspace Scholar: \href{https://scholar.google.com/citations?user=tzu0AM8AAAAJ}{Google Scholar}
\enspace$\vert$\enspace Web: \href{https://datalakes9.github.io/sreelekha.guggilam.github.io/}{Homepage}
\end{center}

% ============================================================================
% RESEARCH INTERESTS
% ============================================================================
\section{Research Interests}
Time-series modeling and databases, Anomaly detection, Spatio-temporal analysis, Machine learning, Quantum computing for data science, Foundation models for time series, Mathematics in Artificial Intelligence

% ============================================================================
%% BEGIN EXPERIENCE
% ============================================================================
\section{Professional Experience}

\textbf{Assistant Professor in Data Science} \hfill August 2024 -- Present\\
Texas A\&M University--Corpus Christi, Department of Mathematics and Statistics\\
Doctoral faculty in Coastal and Marine Sciences Department
\begin{itemize}
    \item Designed graduate course on Advanced Topics in Deep Learning covering generative AI, attention mechanisms, transformers, and model distillation techniques.
    \item Active member of Academics Committee for Data Science PhD Program Planning Committee.
    \item Active member of Graduate Admissions Committee (Mathematics \& Statistics and Coastal \& Marine Sciences).
    \item Active member of Tenure Track Search Committee (Department of Mathematics and Statistics).
    \item Active member of College of Science Curriculum Committee.
\end{itemize}
\vspace{6pt}

\textbf{Research and Development Associate in Machine Learning} \hfill March 2022 -- July 2024\\
Oak Ridge National Laboratory, National Security Sciences Directorate\\
Geo-spatial Science and Human Dynamics Division
\begin{itemize}
    \item DOE Early Career LDRD awardee for research in quantum computing for spatiotemporal modeling. Principal Investigator on developing novel quantum computing approaches for spatiotemporal classification and anomaly detection. Amongst the 8 selected awardees at ORNL in 2023.
    \item Lead on modeling and algorithm development (invention disclosure submitted) for extended Large Deviations Anomaly Detection modeling for indoor pattern of life and human mobility.
    \item Active contributor to multiple Generative AI foundation model projects using multimodal data.
    \item Awarded 4 research grants as Principal Investigator and co-Principal Investigator.
    \item Lead statistician on developing outbreak prediction models using alternate data sources.
\end{itemize}
\vspace{6pt}

\textbf{Research Assistant} \hfill September 2019 -- February 2022\\
University at Buffalo, Institute of Computational Data Science \& Dept.\ of Computer Science and Engineering\\
Mentors: Dr.\ Abani Patra (Tufts University), Dr.\ Varun Chandola (University at Buffalo)
\begin{itemize}
    \item Developed Integrated Clustering and Anomaly Detection (INCAD), a non-parametric unsupervised anomaly detection and clustering framework for streaming multivariate time series using extreme value theory.
    \item Designed Large Deviations Anomaly Detection (LAD), a fast anomaly detection model for high-dimensional multivariate time-series databases. Applied to real-world COVID-19 trend analysis.
    \item Created LAD Inspired Iterative Training (LIIT), a novel low-shot training algorithm for neural networks enabling faster convergence with limited data.
\end{itemize}
\vspace{6pt}

\textbf{Teaching Assistant} \hfill September 2017 -- September 2019\\
University at Buffalo, Institute of Computational Data Science
\begin{itemize}
    \item Led and coordinated a team of teaching assistants for graduate courses in data science.
    \item Collaborated on developing exams and assignments and led recitation sessions.
\end{itemize}
\vspace{6pt}

\textbf{Risk Manager, Global Corporate Portfolios} \hfill November 2015 -- June 2016\\
American Express, World Financial Center, New York, NY
\begin{itemize}
    \item Quantified changes in spending trends in cross-sold clients prior to and post the on-boarding process.
    \item Developed and improved risk margins for corporate clients in Global Corporate Payments.
    \item Identified trends in risk-based industries and states to enable risk control actions on portfolios.
\end{itemize}
%% END EXPERIENCE

% ============================================================================
%% BEGIN EDUCATION
% ============================================================================
\section{Education}

\textbf{Ph.D. in Computational Data Science (CDSE)} \hfill August 2017 -- June 2022\\
University at Buffalo, The State University of New York, Buffalo, NY\\
Dissertation: Non-parametric probabilistic anomaly detection in evolving data: Applications to time series
\vspace{6pt}

\textbf{Masters in Civil Engineering (Transportation Statistics)} \hfill August 2016 -- September 2017\\
University at Buffalo, The State University of New York, Buffalo, NY
\vspace{6pt}

\textbf{Masters in Biostatistics-Bioinformatics} \hfill August 2014 -- September 2015\\
University at Buffalo and Roswell Park Cancer Institute (RPCI), Buffalo, NY\\
Thesis: Statistical assessment of TCGA ovarian cancer sequencing dataset for prognostic utility
\vspace{6pt}

\textbf{Bachelors in Mathematics (B.Math), Honors} \hfill July 2011 -- May 2014\\
Indian Statistical Institute (ISI), Bangalore, India\\
Full student scholarship (merit-based) throughout the completion of the degree
%% END EDUCATION

% ============================================================================
%% BEGIN GRANTS
% ============================================================================
\section{Grant Support}

\textbf{Autonomous Research Institute Fellowship} -- \$82,000 (2 years) \hfill Sep 2025 -- Aug 2027\\
Title: Enhancing Preventive Monitoring through Spatiotemporal Anomaly Detection using Time Series Foundation Models\\
Role: Principal Investigator (100\%)
\vspace{6pt}

\textbf{Port Corpus Christi} -- \$288,000 (1 year) \hfill Apr 2025 -- Dec 2025\\
Title: Vulnerability Index Development -- Phase I -- Data Collection\\
Role: co-Principal Investigator (Share: \$8,000)
\vspace{6pt}

\textbf{DOE-ORNL Early Career Development Research Award} -- \$300,000 (2 years) \hfill Mar 2023 -- Feb 2025\\
Title: Quantum Variational Inference for Anomaly Detection in Spatiotemporal Data\\
Role: Principal Investigator (100\%)
\vspace{6pt}

\textbf{ORNL Laboratory Directors R\&D SEED Funds} -- \$190,000 (2 years) \hfill Feb 2023 -- Jan 2025\\
Title: Artificial Intelligence Models for Land Cover Forecasting\\
Role: Contributor
\vspace{6pt}

\textbf{ORNL Laboratory Directors R\&D Funds} -- \$930,000 (2 years) \hfill Feb 2023 -- Jan 2025\\
Title: Environmental Anomaly Detection for Biopreparedness\\
Role: Co-Principal Investigator
\vspace{6pt}

\textbf{ORNL Laboratory Directors R\&D Funds} -- \$780,000 (2 years) \hfill Oct 2022 -- Sep 2024\\
Title: Pattern of Life for Nuclear Non-proliferation\\
Role: Co-Principal Investigator
%% END GRANTS

% ============================================================================
% PUBLICATIONS (auto-generated from papers.bib)
% ============================================================================
\nocite{*}
\printbibliography[title=Publications]

% ============================================================================
%% BEGIN TEACHING
% ============================================================================
\section{Teaching}

\textbf{DASC-5307} -- Machine Learning in Data Science \hfill Spring 2026\\
Texas A\&M University--Corpus Christi\vspace{4pt}

\textbf{DASC-5390} -- Advanced Topics in Deep Learning \hfill Spring 2025\\
Texas A\&M University--Corpus Christi\vspace{4pt}

\textbf{EAS 501} -- Introduction to Numerical Mathematics for Computing and Data Science \hfill Fall 2017, Fall 2018\\
University at Buffalo (Teaching Assistant)
\vspace{4pt}

\textbf{EAS 502} -- Introduction to Probability Theory for Data Science \hfill Fall 2017, Fall 2018\\
University at Buffalo (Teaching Assistant)
\vspace{4pt}

\textbf{EAS 503} -- Programming and Database Fundamentals for Data Scientists \hfill Fall 2017, Fall 2018\\
University at Buffalo (Teaching Assistant)
\vspace{4pt}

\textbf{EAS 504} -- Applications of Data Science: Industry Overview \hfill Summer 2017, Summer 2018\\
University at Buffalo (Teaching Assistant)
\vspace{4pt}

\textbf{EAS 508} -- Statistical Learning and Data Mining I \hfill Fall 2017, Fall 2018\\
University at Buffalo (Teaching Assistant)
%% END TEACHING

% ============================================================================
%% BEGIN MENTORSHIP
% ============================================================================
\section{Mentorship}

\textbf{Postdoc Mentees}
\begin{itemize}
    \item Madhavi Pagare \hfill 2025--Present
\end{itemize}

\textbf{PhD Advisees}
\begin{itemize}
    \item Anshika Rani (Advisor), Texas A\&M University--Corpus Christi \hfill 2025--Present
    \item Jahnavi Krishna Koda (Committee Member), Texas A\&M University--Corpus Christi \hfill 2025--Present
    \item Nene Coulibaly (Committee Member), Texas A\&M University--Corpus Christi \hfill 2024--Present
\end{itemize}

\textbf{Masters Advisees}
\begin{itemize}
    \item Vigneshwar Lokoji \hfill 2025--Present
    \item Yejin Hwang, Texas A\&M Corpus Christi \hfill 2024--Present
    \item Nikitha Gugulothu (First Appointment: Oak Ridge National Laboratory) \hfill 2022--2023
\end{itemize}

\textbf{Undergraduate Advisees}
\begin{itemize}
    \item Tri Do (First appointment: Ph.D at Oden Institute, UT Austin) \hfill 2023
    \item Mikolaj Jakowski (University of Tennessee, Knoxville) \hfill 2023
\end{itemize}

\textbf{High School Advisees}
\begin{itemize}
    \item Ruhaan Singh (Farragut High School) \hfill 2022--2024
\end{itemize}
%% END MENTORSHIP

% ============================================================================
%% BEGIN SKILLS
% ============================================================================
\section{Technical Skills}

\textbf{Research Focus:} Anomaly detection, Large Deviations Theory, Bayesian Non-parametric Models, Extreme Value Theory, Probabilistic Modeling, Multivariate Time Series, High-Dimensional Streams, Unsupervised Learning
\vspace{4pt}

\textbf{Deep Learning and AI:} Data and Knowledge Distillation, Transformers, Diffusion Models, Attention Mechanisms, Large Language Models, Foundation Models, Generative AI
\vspace{4pt}

\textbf{Programming:} Python (Pandas, NumPy, SciPy, Scikit-learn, PyTorch, TensorFlow, GeoPandas, MPI, OpenMP), C/C++, R, SQL, MATLAB, SAS
\vspace{4pt}

\textbf{Tools and Platforms:} NVIDIA Rapids, Docker, Hadoop, Spark, Git, HPC/Cluster Computing
%% END SKILLS

% ============================================================================
%% BEGIN TALKS
% ============================================================================
\section{Talks and Presentations}

\begin{itemize}
    \item Time-series Modeling -- Opportunities and Challenges with Emerging Methods, University of Texas Rio Grande Valley \hfill 2025
    \item Anomaly Detection in Spatio-temporal Data, ORNL Early Career Development Program \hfill 2023
    \item Quantum Variational Inference for Spatiotemporal Anomaly Detection, Oak Ridge Scientific Advisory Board \hfill 2023
    \item Quantum computing for spatiotemporal analysis, National Security Sciences Week \hfill 2023
    \item Large Deviations for Accelerating Neural Networks Training, Conference on Data Analysis (CoDA), Santa Fe, NM \hfill 2023
    \item Uncertainty, error and anomalies in ML models of remote sensing data, AGU Fall Meeting \hfill 2022
    \item Classifying anomalous members in multivariate time series using large deviations principle, ICCS \hfill 2022
    \item Classifying Collection of Multivariate Time Series Data, ASA UP-STAT Conference, Buffalo, NY \hfill 2022
    \item Identifying anomalous COVID-19 trends using large deviations, Women in STEM Cooperative (WiSC) \hfill 2021
    \item Integrated clustering and anomaly detection (INCAD) for streaming data, ICCS \hfill 2019
    \item Integrating clustering and anomaly detection, CDSE Days, SUNY Buffalo \hfill 2019
    \item National Program on Differential Equations (NPDE) Workshop, IIT Madras \hfill 2013
\end{itemize}
%% END TALKS

% ============================================================================
%% BEGIN AWARDS
% ============================================================================
\section{Awards and Scholarships}

\begin{itemize}
    \item Top Downloaded Author in Wiley Journal of Statistical Analysis and Data Mining \hfill 2023
    \item Oak Ridge National Lab Early Career LDRD Competition Winner \hfill 2023
    \item Runner up: STEM for Everyone -- Women in STEM Cooperative (WiSC) \hfill 2021
    \item Travel Support: UB Navigate Project for women in STEM \hfill 2018
    \item Travel Support: SAMSI -- Model Uncertainty: Mathematical and Statistical (MUMS) \hfill 2018
    \item Travel Support: SAMSI -- Precision Medicine (PMED) \hfill 2018
    \item Honors degree in Mathematics from Indian Statistical Institute \hfill 2014
    \item Full student scholarship (merit-based) at Indian Statistical Institute \hfill 2011--2014
\end{itemize}
%% END AWARDS

% ============================================================================
%% BEGIN SERVICE
% ============================================================================
\section{Professional Service}

\textbf{Organising Committees}
\begin{itemize}
    \item AGU 2023 Fall Meeting Session on Earth System Digital Twins: Prototypes and Federations \hfill 2023
    \item AGU 2022 Fall Meeting Session on AI/ML Assurance: Applications in Geospatial Science \hfill 2022
    \item KDD Workshop on Data-driven Humanitarian Mapping \hfill 2022
    \item CDSE Days, Institute of Computational Data Science, SUNY Buffalo \hfill 2018, 2019
\end{itemize}

\textbf{Referee Service}\\
Data Mining and Knowledge Discovery, Information Systems, Journal of Computational Science, Journal of Hydrology, IEEE Geoscience and Remote Sensing Letters, Frontiers in Earth, Information Sciences, Neural Networks
\vspace{6pt}

\textbf{Panel Judge}
\begin{itemize}
    \item Judge for Health Sciences and Technology Research Symposium \hfill 2025
    \item Judge for the IGNITE Talks \hfill 2023
    \item Judge for SCUDEM VII -- Modeling with Differential Equations \hfill 2022, 2023
    \item Albert Einstein Distinguished Educator Fellowship (AEF) Program \hfill 2022
\end{itemize}
%% END SERVICE

% ============================================================================
%% BEGIN MEMBERSHIPS
% ============================================================================
\section{Professional Memberships}

\begin{itemize}
    \item Delta Omega Honorary Society in Public Health \hfill 2023--Present
    \item Institute of Electrical and Electronics Engineers (IEEE) \hfill 2022--Present
    \item Association for Computing Machinery (ACM) \hfill 2022--Present
    \item American Geophysical Union (AGU) \hfill 2022--Present
    \item Society for Industrial and Applied Mathematics (SIAM) \hfill 2019--Present
    \item Association for Women in Mathematics (AWM) \hfill 2017--2023
    \item Women in STEM Cooperative (WiSC) \hfill 2017--2022
\end{itemize}
%% END MEMBERSHIPS

\end{document}
